\section{Prospective selection and economic sensitivity}\label{sec:decision86}
The best receipt on a completed catalogue and the service chosen before a new workload are different objects. We now connect a calibration decision to independent original-law evaluation and then separate the consequences of policy cost, resource prices and adoption charges. This distinction is needed even when the numerical assignment problem is solved exactly.

\subsection{A selector evaluated on new economic states}
Let $\pi^1,\ldots,\pi^M$ be complete controllers frozen before evaluation. A separate calibration sample selects $\widehat j$. Suppose the evaluation procedure supplies a simultaneous event $\mathcal E$, of probability at least $1-\alpha$, on which every paired cost difference satisfies
\[
 L_{jk}\le J^{\pi^j}-J^{\pi^k}\le U_{jk}.
\]
Put $L_{jj}=U_{jj}=0$. Construction and evaluation error are included in these endpoints; they are not inferred from a sample of midpoint values alone.

\begin{proposition}[Prospective menu regret]\label{prop:selection86}
On $\mathcal E$, the calibration-selected controller satisfies
\begin{equation}\label{eq:regret86}
 0\le J^{\pi^{\widehat j}}-\min_kJ^{\pi^k}
 \le \max_k U_{\widehat j k}.
\end{equation}
Consequently, $\max_k U_{\widehat j k}\le\delta$ establishes noninferiority to the entire frozen menu at a declared margin $\delta$. If $[L_{jk},U_{jk}]\subseteq[-\delta,\delta]$, the pair satisfies cost equivalence at that margin. Strict superiority requires $U_{jk}<0$ and is not implied by equivalence. All statements remain valid after reporting any collection of pairs on the same simultaneous event.
\end{proposition}

The benchmark in \eqref{eq:regret86} is the best expected-cost member of the frozen menu, not the realized winner on a finite evaluation sample and not the original Bellman optimum. The original-optimum account remains the separate construction theorem. Thus a small menu regret does not replace the paper's economic accuracy requirement. The calibration decision in the experiment is made before the evaluation stream is generated. Reconstructing that frozen stream checks its computation; it does not generate another independent statistical experiment.

\subsection{Price and fee regions, not a fitted monetary conclusion}
Suppose $j$ has a declared resource vector $r_j$ and fee $f_j$. Its total economic cost is
\[
 C_j=\kappa J^{\pi^j}+p^{\mathsf T}r_j+f_j,
 \qquad \kappa\ge0,\quad p\ge0.
\]
The conversion $\kappa$ and resource prices $p$ require economic information outside a normalized simulation. They can instead be left in a declared set. Let $\mathcal U$ be a nonempty compact set of jointly admissible policy-cost vectors, incorporating support, paired contrasts and any additional valid restrictions. Define
$D_{jk}=\sup_{v\in\mathcal U}(v_j-v_k)$.

\begin{theorem}[Uniform robust choice region]\label{thm:prices86}
A service $j$ minimizes total cost for every policy-cost vector in $\mathcal U$ at a given $(\kappa,p,f)$ if and only if
\begin{equation}\label{eq:region86}
 \kappa D_{jk}+p^{\mathsf T}(r_j-r_k)+f_j-f_k\le0
 \quad\hbox{for every }k.
\end{equation}
Using certified upper bounds $\overline D_{jk}\ge D_{jk}$ makes these conditions sufficient. If the admissible price--fee set is a compact polytope and the resource vectors are fixed, it suffices to verify these affine inequalities at all its vertices. Boundaries are ties, not strict preference. For any feasible assignment $z$ from a finite assignment menu, the same statement holds after replacing the contrast by
$\sup_{v\in\mathcal U}\langle z-z',v\rangle$
and the resource and fee differences by their assignment aggregates.
\end{theorem}

The result preserves joint cost information. Summing marginal cost upper bounds can be more conservative than optimizing a contrast over $\mathcal U$. A rational dual verifies a proposed support value when $\mathcal U$ is polyhedral. A price region is not an estimated economic gain: it specifies exactly which external valuations support the decision. If $p=f=0$ and $\kappa>0$, an unresolved cost comparison does not become a strict economic preference. Conversely, a strictly negative policy-cost upper bound gives a positive admissible incremental fee without assigning a price to computational work.

For example, the separate evaluation of the frozen investment menu gives a strictly negative upper endpoint for fine minus coarse calibrated Bernstein cost. At $p=0$ and $\kappa=1$, any nonnegative incremental fine-resolution fee below the negative of that endpoint preserves strict preference for the fine controller. The fine calibrated ReLU--Bernstein comparison in that same evaluation remains unresolved in sign, and no strict fee preference is claimed for that pair. Section~\ref{sec:study86} reports the exact-data-derived endpoints and the full simultaneous family.

\subsection{An accuracy design tied to a declared margin}
For independent paired original-law bins, let $[l_i,u_i]$ enclose the true paired conditional cost in bin $i$. Let the paired cost range have width $H_D$. The retained empirical-Bernstein construction gives an interval whose untruncated width is
\begin{equation}\label{eq:width86}
 (\bar u-\bar l)+\sqrt{2s_l^2\lambda/n}
 +\sqrt{2s_u^2\lambda/n}
 +\frac{14H_D\lambda}{3(n-1)}.
\end{equation}
Directed rounding only enlarges the two square-root terms. Its family allocation is stated before the sample is generated. The first term is a numerical enclosure width, the middle terms reflect sampling variability, and the last is a finite-support concentration charge. This is a width decomposition, not a partition of the realized policy loss.

Bin uncertainty, uncertain acquired actions and arithmetic rounding can interact. Unless separate validated counterfactual calculations isolate their contributions, they should be reported jointly in $\bar u-\bar l$, with ambiguity counts and endpoint precision alongside it. The optimizer-class enlargement in Section~\ref{sec:compatible84} is a different term: for two classes it adds $(b_j-a_j)+(b_k-a_k)$ to an attained-policy interval before any valid intersections. It is not another statistical error probability.

For planning, a distribution-free variance bound gives a conservative sample requirement without pilot outcomes. A bounded sample with range $H_D$ has $s^2\le nH_D^2/[4(n-1)]$. Thus the sum of the two square-root terms in \eqref{eq:width86} is at most $H_D\sqrt{2\lambda/(n-1)}$. To make that term and the support charge at most $\delta/4$ each, it suffices that
\[
 n-1\ge\max\{32H_D^2\lambda/\delta^2,
                  56H_D\lambda/(3\delta)\}.
\]
The enclosure width, rounding reserve and any class enlargement still require their own allowance. A narrow interval alone does not guarantee superiority: the true contrast must be separated from the decision boundary. The retained protocol fixes the sample size, precision and normalized margin before evaluation and reports failure to resolve a comparison rather than extending the sample after seeing its sign.
